(1) 特征线方法把 PDE 写成沿轨道的常微分方程 \(du/dt=u_t+\dot x u_x=0\)，选 \(\dot x=c_0+\alpha u\)。因此 u 沿每条特征线保持初值 \(A\sin(k\xi)\)，且
\[
x(\xi,t)=\xi+[c_0+\alpha A\sin(k\xi)]t,\qquad
u(x(\xi,t),t)=A\sin(k\xi).
\]
这是隐式解，只有映射 \(\xi \to x\) 单调时才定义唯一单值的光滑场。

(2) 特征映射 Jacobian
\[
\frac{\partial x}{\partial\xi}=1+\alpha Ak t\cos(k\xi).
\]
首次 Jacobian 为零的时间是
\[
t_s=\frac1{\alpha Ak},
\]
对应 \(\cos(k\xi)=-1\)，可取 \(\xi_s=\pi/k\)。该处 \(\sin(k\xi_s)=0\)，所以首次交叉位置为 \(x_s=\pi/k+c_0t_s\)，再加任意整周期平移。隐函数求导 \(u_x=Ak\cos(k\xi)/[1+\alpha Akt\cos(k\xi)]\) 在该点趋于负无穷，叫梯度灾变。

(3) 方程也可写成守恒形式
\[
u_t+\partial_x\left(c_0u+\frac{\alpha}{2}u^2\right)=0.
\]
对光滑解乘以 u，得到
\[
\partial_t\frac{u^2}{2}+
\partial_x\left(\frac{c_0u^2}{2}+\frac{\alpha u^3}{3}\right)=0.
\]
在长度为整周期的周期边界区域积分，通量两端抵消，因此首次冲击前 \(\int u^2/2\,dx\) 不变；这与无黏性、无耗散的假设一致。

(4) 当 \(\alpha \to 0\) 或 \(A\to 0\)，非线性速度分裂消失，\(t_s\to\infty\)，波形只以速度 \(c_{0}\) 平移。若 t>ts，标签映射会折叠，多条特征线到达同一 x，原有光滑解不再单值；必须引入弱解、Rankine–Hugoniot 条件及熵选择规则确定实际冲击波，不可继续套用交叉前的正则解。
